% =============================================================
%  Section 5 — Valeur perçue par le client
% =============================================================
\section{Valeur perçue par le client}

\begin{pointcle}
  Chaque outil de la stack dispose d'une \textbf{fiche dédiée} (accessible via la colonne « Fiche~→ » du tableau stratégique) détaillant la valeur perçue client, les gains mesurables, et le guide d'installation complet. Cette section synthétise les arguments de vente par profil de décideur.
\end{pointcle}

\subsection{Principe de vente}

\begin{itemize}
  \item Ne jamais citer le nom des outils techniques lors d'un \textbf{rendez-vous commercial initial}
  \item Parler uniquement de \textbf{résultats} : temps gagné, risques évités, image renforcée
  \item Le nom de l'outil n'est pertinent que lors de la phase de \textbf{validation technique}
  \item Chaque fiche individuelle contient le \textit{verbatim type} qu'un client satisfait utiliserait
\end{itemize}

\subsection{Synthèse des arguments de vente}

\begin{table}[htbp]
  \centering
  \caption{Top arguments par type de décideur}
  \renewcommand{\arraystretch}{1.5}
  \begin{tabularx}{\linewidth}{>{\bfseries}p{3cm} X X}
    \toprule
    \rowcolor{tabheader}
    \textcolor{white}{\textbf{Profil décideur}} &
    \textcolor{white}{\textbf{Argument principal}} &
    \textcolor{white}{\textbf{Outil clé à mettre en avant}} \\
    \midrule
    Dirigeant pressé & « Vous gagnez 3–5\,h par semaine dès le premier mois » & n8n + Metabase \\
    \rowcolor{rowgray}
    Dirigeant anxieux sécurité & « Nous vous alertons avant votre IT si quelque chose se passe la nuit » & Wazuh + GoPhish \\
    Directeur commercial & « Votre CRM se met à jour tout seul depuis vos emails » & Odoo + OpenAI API \\
    \rowcolor{rowgray}
    DAF / Contrôleur & « Vos KPIs arrivent automatiquement chaque lundi matin » & Metabase + Grafana \\
    RRH / Responsable qualité & « Vos documents retrouvables en 3 secondes, archivage conforme » & Paperless-ngx \\
    \rowcolor{rowgray}
    DSI / Responsable IT & « Score de sécurité Microsoft en progression, preuves à l'appui » & Secure Score + Wazuh \\
    \bottomrule
  \end{tabularx}
\end{table}

\begin{attention}
  Ne jamais citer le nom des outils techniques lors d'un rendez-vous commercial initial. Parler uniquement de \textbf{résultats} (temps gagné, risques évités, image renforcée). Le nom de l'outil n'est pertinent que lors de la phase de validation technique ou contractualisation.
\end{attention}
